% SPDX-License-Identifier: Apache-2.0
%
% AXI in TxHDL: the tutorial on the bus. Built by //docs:axi. Every
% listing is included from the tree and every printout and figure is
% what the build made by running the examples.
\documentclass[journal,9pt]{IEEEtran}

\newcommand{\listingsize}{\fontsize{6}{7}\selectfont}
% SPDX-License-Identifier: Apache-2.0
%
% The house preamble, shared by every document in this directory. Loaded
% after \documentclass, before \title. Keeping it in one file is what
% stops two articles drifting apart in font, listing style or figure
% style.
% Computer Modern. IEEEtran selects Times (ptm) by default, so the face has
% to be chosen explicitly.
%
% Latin Modern is the Computer Modern variety used here, and the choice is
% forced rather than aesthetic. Under T1 encoding, plain `cmr` has no Type 1
% outlines unless cm-super is installed, and it is not in the pinned
% distribution, so pdflatex falls back to EC bitmaps and every glyph in the
% PDF becomes a Type 3 font. That was measured: the first build produced a
% document with 10 Type 3 fonts and no Type 1. Latin Modern is Computer
% Modern redrawn with a full T1 repertoire and real .pfb outlines, so the
% same design comes out scalable.
\usepackage[T1]{fontenc}
\usepackage[utf8]{inputenc}
\usepackage{lmodern}
% microtype is not in the pinned TeX distribution. emergencystretch alone
% keeps the two-column measure from producing overfull lines.
\setlength{\emergencystretch}{3em}

\usepackage{amsmath}
\usepackage{amssymb}
\usepackage{amsthm}
\usepackage{booktabs}
\usepackage{array}
% enumitem is not in the pinned TeX distribution; the list spacing below
% does what \setlist would have done.
\usepackage{float}
\usepackage{xcolor}
\usepackage{listings}
\usepackage{tikz}
\usetikzlibrary{arrows.meta,positioning,fit,backgrounds,calc}
\usepackage[hidelinks,bookmarks=true]{hyperref}

% Numbered, definition-style box for a rule the reader is meant to apply.
\theoremstyle{definition}
\newtheorem{rulebox}{Rule}
\newtheorem{example}{Example}

% Unnumbered asides.
\theoremstyle{remark}
\newtheorem*{tip}{Tip}
\newtheorem*{pitfall}{Pitfall}

% Tighter lists, in place of enumitem's \setlist.
\setlength{\topsep}{2pt}
\setlength{\itemsep}{1pt}
\setlength{\parsep}{0pt}

\newcommand{\code}[1]{\texttt{#1}}
\newcommand{\term}[1]{\emph{#1}}

% Syntax highlighting. Four roles, distinguishable in colour and still
% distinguishable in greyscale, because an IEEE article gets printed: the
% keyword is the only bold role, the comment is the only italic one, and the
% two remaining roles differ in lightness as well as in hue.
\definecolor{kwblue}{RGB}{20,60,150}
\definecolor{cmtgreen}{RGB}{90,120,90}
\definecolor{strred}{RGB}{150,50,40}
\definecolor{numgray}{gray}{0.45}
\definecolor{framegray}{gray}{0.70}
\definecolor{bgshade}{gray}{0.97}
% A darker fill for figure cells. The listing background has to stay very
% light to keep code readable; a figure cell has to be clearly shaded or the
% distinction it draws is invisible in print.
\definecolor{cellshade}{gray}{0.90}

% The listing size is the document's choice, because it depends on the
% column: a two-column article sets `\listingsize` to 6pt and keeps its
% listings within 70 columns, a one-column document keeps the body size
% and includes source files at 80. Lines are never broken; a line that does not fit
% overflows the frame, which is visible, rather than wrapping, which is
% not.
\providecommand{\listingsize}{\scriptsize}

% `'` is deliberately not a string delimiter: the language uses it for loop
% labels, as Rust does, so treating it as a quote would swallow the rest of
% the line.
\lstdefinestyle{txhdl}{
  basicstyle=\ttfamily\listingsize,
  keywordstyle=\bfseries\color{kwblue},
  commentstyle=\itshape\color{cmtgreen},
  stringstyle=\color{strred},
  numberstyle=\tiny\color{numgray},
  morekeywords={mod,use,pub,const,type,struct,enum,trait,impl,fn,async,let,mut,
    match,if,else,loop,while,for,in,break,continue,return,as,await,
    module,bus,channel,transaction,protocol,pipeline,flow,seq,config,
    port,stage,pipe,instance,connect,bind,tag,domain,blackbox,foreign,
    property,role,signal,handler,true,false,
    sequence,interface,modport,implements,package,import,var,wire,func,proc,
    case,default,next,fork,branch,join,state,fsm},
  morecomment=[l]{//},
  morestring=[b]",
  showstringspaces=false,
  columns=fullflexible,
  keepspaces=true,
  breaklines=false,
  % Framed, so a listing is bounded rather than floating in the column.
  frame=single,
  rulecolor=\color{framegray},
  framesep=3pt,
  backgroundcolor=\color{bgshade},
  xleftmargin=3pt,
  xrightmargin=3pt,
  aboveskip=6pt,
  belowskip=6pt,
  % Every listing is numbered and captioned, so it can be referred to by
  % number. `caption` without `float` numbers the listing and leaves it
  % where it was written, which is what a listing in running prose needs.
  captionpos=b,
  abovecaptionskip=2pt,
  belowcaptionskip=6pt,
}

% Figures drawn as text. Framed the same way, and with the highlighting
% switched off, because a box-and-arrow diagram is not code and half its
% words turning blue reads as an accident. Setting a style adds keys rather
% than clearing the ones already set, so the three roles are neutralised by
% hand rather than by leaving them out.
\lstdefinestyle{ascii}{
  basicstyle=\ttfamily\listingsize,
  keywordstyle=\ttfamily,
  commentstyle=\ttfamily,
  stringstyle=\ttfamily,
  columns=fullflexible,
  keepspaces=true,
  frame=single,
  rulecolor=\color{framegray},
  framesep=3pt,
  backgroundcolor=\color{bgshade},
  xleftmargin=3pt,
  xrightmargin=3pt,
  aboveskip=6pt,
  belowskip=6pt,
}
\lstset{style=txhdl}

% House TikZ styles. `shorten` lives on the arrow styles rather than on each
% arrow, so every arrow in the document gets the gap, including ones added
% later. TikZ clips a path at a node's border, so without it an arrowhead
% butts against the box with no gap at all. Keep `node distance` well above
% twice the shorten, or the arrow shrinks to nothing.
\tikzset{
  box/.style={draw,rounded corners=2pt,align=center,inner sep=4pt,
              font=\footnotesize},
  boxf/.style={box,fill=bgshade},
  lbl/.style={font=\scriptsize\itshape,align=center},
  fl/.style={-{Stealth[length=4pt]},shorten >=2.5pt,shorten <=2.5pt},
  fld/.style={fl,dashed},
}



\InputIfFileExists{buildstamp.tex}{}{}
\providecommand{\buildstamp}{Draft build (outside the Bazel flow).}

% A range of the library's source, by its begin/end markers.
\newcommand{\axipart}[3]{%
\lstinputlisting[rangebeginprefix=//\ begin\{,rangebeginsuffix=\},%
rangeendprefix=//\ end\{,rangeendsuffix=\},includerangemarker=false,%
linerange=#1-#1,caption={#2},label={#3}]{lib/parts/src/bus/axi.rs}}

\title{AXI in TxHDL}

\author{Filip Filmar and Dragiša Janković%
\thanks{Both authors are independent. Filip Filmar is the
corresponding author, at f@hdlfactory.com; Dragiša Janković is
reached at linkedin.com/in/dragisa-jankovic.}%
\thanks{This is the tutorial on one bus of the language's library,
for a reader who has met TxHDL through the step-by-step
tutorial~\cite{tutorial} and wants to put an AXI link in a design,
or to read how the link is made. Every listing is included from the
project repository \code{HDL/txhdl} at build time and every printout
and figure is what the build produced by running the examples. The
reservation station has a tutorial of its own~\cite{station}, the
runtime is in the runtime document~\cite{runtime}, and the
cover~\cite{cover} lists every document. The AXI examples are here
and not in the examples document~\cite{examples}, because a bus is a
subsystem and not one more construct.}%
\thanks{The authors used a large language model, Claude, as an
assistant in exploring the concepts and in writing and constructing
the documents and the programs; every commit of the repository says
so and carries the prompts that led to it, verbatim.}%
\thanks{\buildstamp}}

\markboth{AXI in TxHDL}{Filmar and Janković: AXI in TxHDL}

\begin{document}
\maketitle

\begin{abstract}
A channel carries one transaction in one direction. AXI is five
channels that only mean something together, and a client written
against them directly must gather a write address with its data
beats, count them, watch \code{last}, allocate identifiers and match
them back. TxHDL provides the link instead: \code{axi()} makes every
channel of an AXI4 link at once and hands back two ends, the way
\code{chan()} hands back a sender and a receiver. A host client
issues a burst and awaits its answer; a peripheral client accepts a
whole decoded transaction and answers it, from whatever process it
likes and in whatever order its work finished. Neither names a
channel. Between the two ends are the two units that do the
protocol, lowered to Verilog and to VHDL and simulated against the
examples' runs, with the five AXI channels between them, which is
where a router will go.
\end{abstract}

\section{What a Link Is}
\label{sec:a-what}

AXI4 is five channels. Write address, write data and write response
go \code{aw}, \code{w}, \code{b}; read address and read data go
\code{ar} and \code{r}. A write is an address beat, one or more data
beats with the last of them marked, and a response beat. A read is
an address beat and one or more data beats, again with the last
marked. Every beat carries an identifier, and beats of different
identifiers may be interleaved and answered in any order, which is
what lets several transactions be in flight at once.

None of that is a client's work, and in TxHDL none of it is written
by one. \code{axi()} makes the five channels, and the channels
between each end and the unit that speaks the protocol for it, and
returns a \code{Link}: a \code{Host} end, a \code{Per} end, and the
ports of the two units. A design joins the two units as it joins any
others, and its clients hold the two ends.

\begin{figure}[htbp]
\centering
\begin{lstlisting}[style=ascii,frame=none,numbers=none]
  host client        issue / beat / grant / done / release
  AxiHost            a unit, lowered
     aw ar w b r     the AXI4 channels, where a router goes
  AxiPer             a unit, lowered
  peripheral client  request / write data / answer / read beat
\end{lstlisting}
\caption{The layers of a link. The middle is real AXI; everything
above and below it is transactions.}
\label{fig:a-layers}
\end{figure}

\section{The Client}
\label{sec:a-client}

A host client has two verbs. \code{read} takes an address and a
count of beats; \code{write} takes an address and the beats
themselves. Each returns when the burst has gone out, with a handle
that is the answer to come. Several handles may be outstanding, and
they may be awaited in any order.

\axipart{hostapi}{The host end, from
\code{lib/parts/src/bus/axi.rs}: the whole of what a client that
issues bursts uses.}{lst:a-hostapi}

The burst's length is the slice the client passed, so it cannot
disagree with the beats that follow it; the last beat is marked as
the last by \code{write} and not by the client; and the identifier is
allocated by the hardware and told back, so a client never writes
one. What a client writes is the address and the data.

A peripheral client has one verb. \code{accept} waits for a whole
AXI transaction and returns it decoded: a read, with the count of
beats it wants, or a write, with every beat of its data already
gathered. The transaction answers itself, with \code{ok}, \code{err}
or \code{data}.

\axipart{accept}{The peripheral end: one call, and what comes back
is a whole transaction.}{lst:a-accept}

A read is ready as soon as its address beat arrives. A write is
ready when its last data beat has, because AXI4 puts no identifier
on the write data channel, so the beats belong to the oldest address
beat and there is nothing to gather them by but order. One accept
gathers at a time, which is what lets several processes share one
end without taking each other's beats.

\section{Accepting Many}
\label{sec:a-many}

A peripheral that does real work does not answer one transaction
before accepting the next. It accepts as they arrive, hands each to
whatever does the work, and answers each as its data becomes
available. An accepted transaction carries its own identifier and a
handle on the end's answer port, so it may be moved to another
process and answered there, long after and out of order.

\code{Open} holds the ones in between, by identifier, and
\code{serve} is the whole shape as one call: as many processes as
the design wants slots, each accepting when it is free and answering
when its work finished.

\axipart{serve}{\code{serve}: \code{pipeline::drive} for a bus. The
client writes the body and nothing of the protocol.}{lst:a-serve}

\code{ex\_axi\_serve} writes the same peripheral twice, once by hand
and once with \code{serve}, and asserts that the two answer alike.
By hand it is three processes: one accepts and forwards a ticket to
a pool of workers, the workers take their own time, and a third
answers whatever finished, taking the transaction back out of
\code{Open}. That third process never called \code{accept}.

\begin{figure}[htbp]
\centering
\input{axi_serve_timing.tex}
\caption{Ten one-word reads over four identifiers, three in flight.
An identifier goes back on \code{release} only once the client has
taken that burst's answer, and \code{busy} shows which are out.}
\label{fig:a-serve}
\end{figure}

\lstinputlisting[caption={\code{ex\_axi\_serve.rs}. Run with
\code{bazel run //lib/examples:ex\_axi\_serve}.},label={lst:a-serve-ex}]{lib/examples/ex_axi_serve.rs}

\lstinputlisting[style=ascii,caption={What \code{ex\_axi\_serve}
printed when the build ran it: the two peripherals, each answering
in the order its work finished.},label={lst:a-serve-out}]{out_axi_serve.txt}

\section{The Beats}
\label{sec:a-beats}

The values on the wire are AXI4's, in full: the address phase
carries \code{id}, \code{addr}, \code{len}, \code{size},
\code{burst}, \code{lock}, \code{cache}, \code{prot}, \code{qos} and
\code{region}, and it is one struct, because \code{AW} and \code{AR}
are the same fields under different names.

\axipart{beats}{The five channels' beats.}{lst:a-beats}

The widths are stated rather than computed, as they are everywhere
in this library: \code{A} is the address width, \code{D} the data
width, \code{S} the strobe width, which is \code{D/8}, \code{I} the
identifier width and \code{NIDS} the count of identifiers, which is
\code{1 << I}. \code{U<\{D / 8\}>} would need nightly Rust, and the
prototype pays that in stated parameters rather than in a nightly
toolchain, exactly as \code{Fifo} states its depth and
\code{Station} its line count.

\section{The Trackers}
\label{sec:a-trackers}

Two units do the protocol, and both are written in the lowered
subset, so each is a netlist as well as a simulation.

\code{AxiHost} turns an issue into an address beat on whichever
channel it belongs to, hands out the identifiers, passes the write
beats out, and reports back what returns. The identifiers go round
robin, one per burst.

\axipart{hostunit}{\code{AxiHost}: the host side, as
hardware.}{lst:a-hostunit}

An identifier is outstanding from its address beat until the client
releases it, and the client's end releases it when it has taken that
burst's answer. That is one cycle later than freeing it on the
response would be, and it is the difference between correct and
almost correct: an identifier freed on its response could be handed
out again while the answer to the burst before still stood
uncollected, and the two answers would be gathered into one. A burst
whose turn falls on an identifier still outstanding waits, which is
the backpressure that bounds how many transactions are in flight.

\code{AxiPer} merges the two address channels into one stream of
requests, hands the write beats to the client, and turns the
client's answers back into response beats. The two address channels
take turns, so neither starves the other however long one of them
streams.

\axipart{perunit}{\code{AxiPer}: the peripheral side.}{lst:a-perunit}

\section{Supplying Your Own}
\label{sec:a-own}

Neither unit is privileged. \code{axi()} returns their ports as
named types, \code{HostIn}, \code{HostOut}, \code{PerIn} and
\code{PerOut}, so a design writes a unit of its own against the same
ports and joins that instead. A peripheral that is real hardware may
also ignore \code{AxiPer} altogether and read the five AXI channels
directly, which is what a memory controller does.

\section{A Client That Is Hardware}
\label{sec:a-hardware}

\code{Host} and \code{Per} are simulation-side objects. \code{Host}
keeps an inbox of answers and hands out futures; \code{Per} gathers a
write's beats into a vector; \code{Open} is a table keyed by
identifier and \code{serve} starts as many processes as it is asked
for. None of that lowers, and none of it is meant to: they are what a
testbench, or a model, or the software side of a design holds.

Underneath them are channels, and a client written as a unit holds
those channel ends itself. On the host side they are the issue
channel, the write beats and the release going out, and the grant,
the write responses and the read beats coming back; on the peripheral
side, the requests and write beats coming in and the answers and read
beats going out. \code{axi\_units} makes a link and hands those out
under the names \code{HostClient} and \code{PerClient}, which is the
same link \code{axi} makes, with the two ends unwrapped.

\begin{lstlisting}[style=ascii,frame=none,numbers=none]
  axi()        -> Host / Per     a client that is a testbench
  axi_units()  -> HostClient /   a client that is a unit, and
                  PerClient      lowers with the rest of a design
\end{lstlisting}

A rasteriser written this way is the GPU's~\cite{gpu}, and it is the
first client of this link that is hardware. Nothing about the link
changes between the two: the same two trackers, the same five AXI
channels, the same identifiers.

\section{Peer to Peer}
\label{sec:a-peer}

\code{ex\_axi} is one host and one peripheral on one link. The host
issues three bursts, two reads and a write, without waiting for any
of them, so all three are in flight with an identifier each. The
peripheral serves them with three slots and a worker whose time
depends on the address, so the work finishes in an order of its own.

\begin{figure}[htbp]
\centering
\input{axi_timing.tex}
\caption{The peer-to-peer run: three bursts issued, the identifiers
granted, the address beats out on \code{ar} and \code{aw}, the write
beats with \code{w\_last}, and the answers back in an order that is
not the order they were issued.}
\label{fig:a-peer}
\end{figure}

\lstinputlisting[caption={\code{ex\_axi.rs}. Run with
\code{bazel run //lib/examples:ex\_axi}.},label={lst:a-peer-ex}]{lib/examples/ex_axi.rs}

\lstinputlisting[style=ascii,caption={What \code{ex\_axi} printed
when the build ran it: the three bursts, the order they were
answered in, and then the lowering, two modules.},label={lst:a-peer-out}]{out_axi.txt}

The client awaits its three handles in the order it issued them and
still gets each burst its own answer, although the answers arrived
in another order. An answer that arrives early waits in the host
end's inbox until the handle it belongs to is awaited; the drain
that puts it there is done by whichever handle is being polled, so
no process has to be joined for it and no answer is ever stuck
behind another at the head of a channel.

\section{The Netlist}
\label{sec:a-netlist}

The end of Listing~\ref{lst:a-peer-out} is the Verilog of the two
trackers, as the lowering wrote it. Every channel port becomes
\code{\_data}, \code{\_valid} and \code{\_ready}, so \code{AxiPer}
is a module with \code{aw\_data}, \code{aw\_valid},
\code{aw\_ready}, \code{ar\_data} and the rest, which is the shape a
synthesis tool expects of an AXI subordinate, with each channel's
fields packed into one word rather than split into \code{awaddr},
\code{awlen} and the others. Splitting them is a renaming of slices
and nothing more.

The identifiers are a register and a mask of four bits; the taking
turns is one register of one bit; every wire of the step is named
for what it is: \code{free}, \code{go}, \code{wr\_go},
\code{rel\_go}, \code{take\_ar}, \code{take\_aw}. Both netlists are
generated by the run of \code{ex\_axi} and simulated against that
run by the build, as VHDL under \code{nvc} and as Verilog under
Verilator, at their ports, every time.

One name is worth the warning. A register called \code{next} lowers
to a VHDL signal called \code{next}, which VHDL reserves, and the
netlist then does not analyse; the register here is called
\code{turn} for that reason. The lowering does not yet escape the
reserved words of its targets.

\section{What Routing Attaches To}
\label{sec:a-routing}

The five AXI channels between \code{AxiHost} and \code{AxiPer} are
ordinary channels of ordinary transactions. A router is a unit with
an \code{aw}, \code{ar} and \code{w} in and several out, and several
\code{b} and \code{r} in and one out, choosing by address on the way
down and by identifier on the way back, which is what
\code{cpu/vreteno/src/router.rs} already does for that core's
simpler bus. Nothing in the two ends changes when one is inserted:
the client verbs are the same, and the identifiers, which the
clients never see, are what the answers come home by.

\begin{thebibliography}{9}

\bibitem{tutorial}
F.~Filmar and D.~Janković, ``TxHDL, step by step,'' project
repository \code{HDL/txhdl}, \code{//docs:tutorial}, 2026.

\bibitem{runtime}
F.~Filmar and D.~Janković, ``The TxHDL runtime,'' project repository
\code{HDL/txhdl}, \code{//docs:runtime}, 2026.

\bibitem{examples}
F.~Filmar and D.~Janković, ``TxHDL by example,'' project repository
\code{HDL/txhdl}, \code{//docs:examples}, 2026.

\bibitem{station}
F.~Filmar and D.~Janković, ``A reservation station in TxHDL,''
project repository \code{HDL/txhdl}, \code{//docs:station}, 2026.

\bibitem{gpu}
F.~Filmar and D.~Janković, ``A minimal GPU in TxHDL,'' project
repository \code{HDL/txhdl}, \code{//docs:gpu}, 2026.

\bibitem{cover}
F.~Filmar and D.~Janković, ``TxHDL documents,'' project repository
\code{HDL/txhdl}, \code{//docs:cover}, 2026.

\end{thebibliography}

\end{document}
